%
% reversi_280_momento_evolutivo.tex
%
% This file is part of the reversi program
% http://github.com/rcrr/reversi-v2
%
% Author Roberto Corradini mailto:rob_corradini@yahoo.it
% Copyright 2026 Roberto Corradini. All rights reserved.
%
% License
%
% This program is free software; you can redistribute it and/or modify it
% under the terms of the GNU General Public License as published by the
% Free Software Foundation; either version 3, or (at your option) any
% later version.
%
% This program is distributed in the hope that it will be useful,
% but WITHOUT ANY WARRANTY; without even the implied warranty of
% MERCHANTABILITY or FITNESS FOR A PARTICULAR PURPOSE. See the
% GNU General Public License for more details.
%
% You should have received a copy of the GNU General Public License
% along with this program; if not, write to the Free Software
% Foundation, Inc., 59 Temple Place - Suite 330, Boston, MA  02111-1307, USA
% or visit the site <http://www.gnu.org/licenses/>.
%

%-----------------------------------------------------------------------------
% Capitolo: Il momento evolutivo della partita.
%
% Prima fase: il modello 0108 (configurazione di riferimento) e' eseguito a ogni
% valore di empty_count da 24 a 4, a passo 1: 21 regressioni. Per
% empty_count 25-60 non esistono dati nel database REGAB; la serie non si
% estende a quell'intervallo. La RMSE di validazione e' in fascia stabile
% (6.96-7.05) da 11 a 24 e decresce a ogni passo da 10 a 4 (6.88 -> 5.65);
% i pesi attivi calano da 80.457 a 66.962.
%-----------------------------------------------------------------------------

\chapter{Il momento evolutivo della partita}

\section{La RMSE di validazione per empty count}

Il modello 0108 (configurazione di riferimento: dieci pattern e la
mobilità OptMob, addestrato sui batch 3, 9 e 12, validato sul batch
6) è eseguito a ogni valore di \texttt{empty\_count} da 24 a 4, a
passo 1: 21 regressioni, una per valore. Per \texttt{empty\_count}
compresi tra 25 e 60 non esistono dati nel database REGAB; la serie
non si estende a quell'intervallo. La
tabella~\ref{tab:ec-rmse} riporta, a ogni valore, la RMSE di
validazione, i pesi attivi $W$ e le iterazioni di ottimizzazione.

\begin{table}[ht]
\centering
\begin{tabular}{rrrr}
\toprule
empty count & RMSE validazione & pesi attivi $W$ & iterazioni \\
\midrule
24 & 6.99 & 80\,457 & 774 \\
23 & 7.00 & 80\,696 & 711 \\
22 & 6.98 & 80\,974 & 718 \\
21 & 7.02 & 81\,079 & 715 \\
20 & 7.01 & 81\,269 & 687 \\
19 & 7.04 & 81\,372 & 714 \\
18 & 7.00 & 81\,547 & 702 \\
17 & 7.05 & 81\,620 & 703 \\
16 & 7.00 & 81\,714 & 683 \\
15 & 7.02 & 81\,775 & 738 \\
14 & 7.01 & 81\,816 & 721 \\
13 & 7.02 & 81\,797 & 710 \\
12 & 6.97 & 81\,822 & 743 \\
11 & 6.96 & 81\,787 & 716 \\
10 & 6.88 & 81\,739 & 712 \\
9  & 6.84 & 81\,627 & 778 \\
8  & 6.75 & 81\,502 & 770 \\
7  & 6.62 & 81\,063 & 840 \\
6  & 6.39 & 79\,785 & 886 \\
5  & 6.13 & 75\,932 & 941 \\
4  & 5.65 & 66\,962 & 941 \\
\bottomrule
\end{tabular}
\caption{La RMSE di validazione, i pesi attivi e le iterazioni al
variare di \texttt{empty\_count} (il modello 0108, configurazione di
riferimento; 21 regressioni indipendenti).}
\label{tab:ec-rmse}
\end{table}

Per \texttt{empty\_count} da 11 a 24 la RMSE di validazione resta
nell'intervallo 6.96--7.05 (ampiezza 0.09); il valore a 20 è 7.01.
Per \texttt{empty\_count} da 4 a 10 la RMSE decresce a ogni passo:
da 6.88 a 10 a 5.65 a 4. I pesi attivi diminuiscono al diminuire
di \texttt{empty\_count}: da 80\,457 a 24 a 66\,962 a 4. Le iterazioni di
ottimizzazione sono comprese tra 683 e 941, in crescita al
diminuire di \texttt{empty\_count}.
